% status: ZERO
% Printed Chapter 34 / stable source ch17. Rules: AUTHORING.md.
\chapter{Prefix Caching and RadixAttention}\label{ch:prefix-caching}

\begin{ChapterIntent}
Paged KV memory gave cached state physical identity, generations, reference counts,
immutable sharing, and safe retirement. Prefix caching adds a different problem:
recognition. Given a new request, can the engine prove that some already-computed KV
state is exactly the state this request would compute for its prefix? This chapter
derives that proof obligation, implements hash-chain and radix-tree indexes, separates
logical cache identity from physical block ownership, makes tenant isolation part of
the key space, and treats eviction as an economic decision between retained bytes and
avoided prefill. RadixAttention is presented correctly as a runtime reuse and
scheduling architecture, not as an attention kernel.
\end{ChapterIntent}

\OpeningSection{The computation you should not pay for twice}

The measurement record used by this edition contains a useful experiment on an Apple
M1 running llama.cpp with Llama~3.2~3B. A roughly 3,000-token document followed by a
question required about 11.5 seconds before the first token on a cold request. A
second request over the same document but a different suffix arrived in about
0.16 seconds after the server reported reusing roughly 3,004 cached tokens and
processing only the short new suffix. These are measured values from
\texttt{research/measurements/2026-09-24-m1-part3.md}; they are not universal
performance constants.

The dramatic speedup is easy to misunderstand. The engine did not discover that two
pieces of prose were semantically similar. It proved a much narrower fact: a long
causal prefix of the second request had already executed under a compatible inference
configuration, and the resulting KV state was still available. Change a token near
the front and exact reuse stops at that divergence. Move the same document later in
the prompt and its KV is generally different because its preceding context and
positions changed.

This makes prefix caching a systems problem with a crisp contract. The index may
\emph{propose} a reusable prefix. The engine may consume it only if three independent
conditions hold: the computation identity matches, the reuse boundary is legal, and
the physical state named by the index is still live and protected from reuse. A fast
hash lookup is not itself that proof.

\Foundations{A cache hit is a claim about computation, not text}{
For a causal transformer, the KV state at logical position $p$ is determined by the
model computation over positions $0\ldots p$. Equal rendered strings are insufficient.
The token IDs, model and weight revision, adapter, positional convention, relevant
multimodal inputs, KV representation, and any other state that changes the forward
pass belong to the computation namespace. Sampling temperature does not normally
belong there because it acts after logits are produced. Chapter~\ref{ch:paged-kv}
supplies the physical contract: immutable shared blocks may be referenced by multiple
requests, mutable tails require ownership or copy-on-write, and generations prevent a
stale logical reference from silently naming a recycled physical block.
}

\NapkinSection{Napkin math: when are retained bytes worth it?}

Let $L$ be the number of prefix tokens a request can reuse, $P$ the model parameter
count, $m_{kv}$ the retained KV bytes per token across all layers, and $h$ the
probability that the retained prefix is reused before eviction. Using the familiar
first-order estimate of about $2P$ FLOPs per prefill token, the expected compute
avoided by retaining the prefix is
\begin{equation}
E[F_{\rm saved}] \approx h\,2PL .
\end{equation}
The retained memory is
\begin{equation}
M_{\rm kept}=L\,m_{kv}.
\end{equation}
Their ratio is
\begin{equation}
\rho = \frac{E[F_{\rm saved}]}{M_{\rm kept}}
      \approx h\frac{2P}{m_{kv}}
      \quad [\mathrm{FLOP/byte}].
\label{eq:prefix-value}
\end{equation}
This is a derived exchange rate, not a complete eviction policy. It omits the time
until reuse, memory pressure, transfer cost between tiers, recomputation overlap,
scheduler interference, and the fact that HBM capacity may be worth more during a
burst than while the server is idle.

A more operational score for cached object $j$ is
\begin{equation}
V_j = \frac{\hat p_j C_j - X_j}{M_j\,\Delta t_j},
\label{eq:cache-value}
\end{equation}
where $\hat p_j$ is estimated reuse probability, $C_j$ is recomputation cost avoided,
$X_j$ is expected restore/management cost, $M_j$ is resident bytes, and $\Delta t_j$
is expected time until reuse. Equation~\ref{eq:cache-value} is illustrative: an
industrial policy may approximate these terms with recency, frequency, prefix depth,
tenant priority, or measured prefill time.

\section{The reuse theorem the runtime actually needs}

Write the computation namespace as
\[
\mathcal N=(W,\tau,A,\Pi,Q,\mu,\kappa),
\]
where $W$ identifies model weights, $\tau$ the tokenizer contract, $A$ the active
adapter set, $\Pi$ positional/attention semantics, $Q$ the KV representation and
quantization contract, $\mu$ relevant multimodal content identity, and $\kappa$ the
security namespace. The exact tuple is engine-specific; the engineering rule is not.

For token sequence $x_{0:n}$, let $K_{\mathcal N}(x_{0:n})$ denote the deterministic
KV state produced for that prefix under namespace $\mathcal N$. Exact prefix reuse
from cached request $a$ into request $b$ through position $r-1$ is valid when
\[
\mathcal N_a=\mathcal N_b,\qquad
x^a_{0:r}=x^b_{0:r},
\]
the cache contains a legal checkpoint at $r$, and every physical object referenced by
that checkpoint is still valid. This is the central invariant of the chapter.

The security namespace $\kappa$ deserves emphasis. If two tenants use identical
tokens but policy forbids cross-tenant reuse, they must inhabit different cache key
spaces even though their mathematical transformer state would match. Correctness and
authorization are separate dimensions of validity.

\section{Why a middle substring is not a prefix}

Suppose requests are $A\|D$ and $B\|D$, where $D$ is the same retrieved document.
The document bytes match, but the hidden states and hence keys and values for $D$
depend on $A$ versus $B$. Reusing $D$ merely because its token subsequence matches
would substitute a state computed under the wrong causal history.

This explains a practical prompt-layout rule: place stable, frequently reused
material before volatile material when application semantics permit. A timestamp,
nonce, request ID, or reordered tool schema near the beginning destroys the long
common prefix even when everything after it is identical. Prefix caching rewards
structural prompt stability, not textual redundancy in arbitrary locations.

\section{Hash-chained blocks: content addressing with ancestry}

Paged KV naturally suggests block-level recognition. Let $t_i$ be the token IDs in
logical block $i$, $e_i$ the extra computation/security dependencies that first
become relevant there, and $H_{i-1}$ the identity of the preceding prefix. Define
\begin{equation}
H_i = H(H_{i-1}, t_i, e_i).
\label{eq:hash-chain}
\end{equation}
Because ancestry participates in the key, identical block tokens reached through
different prefixes do not alias.

The index maps $H_i$ to one or more physical block handles. A lookup walks left to
right until the first unavailable legal checkpoint. The index is metadata; the block
pool remains the authority on physical liveness. A candidate handle carrying
$(id,generation)$ must still validate before the scheduler pins it.

Current vLLM uses this family of design: parent hash, block tokens, and extra hashes
such as LoRA identity, multimodal hashes, and cache salt contribute to cache identity.
Its current documentation also distinguishes the prefix matching unit from physical
KV block size for configurations where a finer legal match boundary is supported.
Those are current implementation facts and must be pinned to an exact upstream commit
when this chapter is integrated.

\section{Hash collisions are a correctness boundary}

A non-cryptographic hash is not equality. If an implementation maps a digest directly
to KV without collision verification, a collision can turn a performance mechanism
into silent wrong-model-state reuse. There are three defensible patterns: use a
cryptographically strong digest with an explicit risk budget; retain sufficient
canonical key material to verify equality after the hash lookup; or combine both.

The lab takes the pedagogically strongest route: a deliberately weak hash narrows the
candidate set, but exact namespace and token-prefix equality are checked before a hit
is accepted. A broken variant skips verification; the tests manufacture a collision
and demonstrate false reuse.

\EngineFigure{ch17-hash-proof}{A hash lookup proposes a candidate; namespace equality, exact prefix equality, legal checkpoint, and live generation collectively authorize reuse.}{fig:prefix-hash-proof}

\section{Radix trees: represent the prefix relation directly}

A radix tree is a compressed trie. Edges carry token spans rather than one token each.
A common system prompt occupies a path near the root; tool schemas, few-shot examples,
retrieved contexts, and conversation turns create branches. Matching a new request
walks the tree until the first token divergence. Insertion may split an edge at the
longest common prefix.

This representation exposes workload structure that a flat hash table hides. The
runtime can ask not only ``is block hash $h$ resident?'' but ``which active or cached
request shares the deepest prefix with this arrival?'' That information can influence
scheduling and routing.

SGLang's RadixAttention is built around this observation. The name is easy to
misread: it is not a replacement for scaled dot-product attention. It is a runtime
architecture for reusing KV state across shared prefixes and coordinating execution
around the radix structure. The original SGLang paper reports up to 6.4$\times$
throughput on its evaluated structured workloads; that is an authors' measurement,
not a universal multiplier.

\EngineFigure{ch17-radix-tree}{A compressed radix tree makes common causal ancestry explicit. Active paths are pinned; unpinned leaves are natural eviction candidates.}{fig:prefix-radix}

\section{Hash table or radix tree?}

The two structures solve overlapping but not identical problems. Hash chains are
compact and align naturally with fixed-size paged blocks. They provide direct lookup
for a known checkpoint key and distribute well. A radix tree preserves prefix
topology, supports longest-prefix matching directly, and exposes locality useful to a
scheduler.

Production systems can combine ideas. Physical storage can remain paged while a radix
index names spans; a radix node can point to block handles; hashes can protect or
accelerate edge identity. The false dichotomy ``hash caching versus RadixAttention''
is therefore less useful than asking which layer owns identity, topology, storage,
pinning, and eviction.

\section{The partial-block boundary}

Chapter~\ref{ch:paged-kv} established that a shared mutable tail is dangerous. Prefix
caching inherits the rule. The simplest design publishes only full immutable blocks.
If physical blocks hold $B$ tokens, a 19-token common prefix with $B=16$ reuses 16
tokens and recomputes three.

More advanced engines may support a match granularity finer than physical storage.
That does not abolish the proof obligation: the resumed state must correspond to a
checkpoint the model/cache architecture can legally restore. Hybrid attention,
sliding-window state, recurrent state-space layers, and multimodal models can make
``checkpoint at every token'' false or expensive. Logical match granularity and
physical allocation granularity must therefore be modeled separately.

\section{From candidate to pinned lease}

A robust hit path is transactional:
\begin{enumerate}
\item tokenize and construct the computation/security namespace;
\item find the longest candidate prefix in the hash/radix index;
\item verify exact identity and a legal checkpoint;
\item validate every physical handle generation;
\item increment references/pins, removing newly active blocks from eviction eligibility;
\item attach the handles to the request block table;
\item only then report the tokens as computed to the scheduler.
\end{enumerate}

If any validation fails, the request falls back to a shorter valid prefix or cold
prefill. The index must never be able to resurrect a block merely because an old
digest remains in a map.

\section{Eviction is admission control for future work}

Cached KV competes with active KV for the same scarce memory. Keeping everything
maximizes theoretical hit rate and eventually prevents new work from being admitted.
Evicting everything maximizes free capacity and destroys reuse. The allocator and
cache therefore need a boundary: active or pinned blocks are not eviction candidates;
unreferenced cached blocks are.

LRU is attractive because it approximates temporal locality with tiny metadata.
Radix trees add a structural constraint: removing an interior representation may
affect descendants, so implementations often reason about evictable leaves or
account for shared ancestry. Cost-aware policies can prefer a long expensive prefix
over a short cheap one even when the short prefix is more recent.

A useful metric is not hit count but \emph{prefill work avoided}. Ten hits on
16-token prefixes can be less valuable than one hit on a 4,000-token system prompt.
Record hit tokens, avoided prefill time, bytes retained, eviction churn, and the
effect on admission failures.

\section{A worked cache-pressure example}

Assume 12 physical blocks are available and active requests pin seven. Five blocks
remain as the cache budget. Prefix $A$ occupies three blocks and costs 30 ms per block
to recompute; prefix $B$ occupies two and costs 8 ms per block. If $A$ has estimated
reuse probability 0.4 and $B$ 0.8, their expected avoided compute is respectively
$0.4(90)=36$ ms and $0.8(16)=12.8$ ms. Per block, $A$ still has higher expected
compute value: 12 ms/block versus 6.4 ms/block. An LRU policy might make the opposite
choice if $B$ was touched more recently.

Neither policy is ``correct'' without a service objective. If retaining $A$ causes a
high-priority request to queue for lack of blocks, the opportunity cost can dominate
both estimates. Cache policy is subordinate to scheduler policy.

\section{Cache-aware scheduling and routing}

Once prefix locality is observable, the scheduler can exploit it. Among otherwise
similar queued requests, executing one whose prefix is already resident can reduce
prefill tokens and TTFT. Across replicas, routing a request to a worker with a warm
prefix avoids paying the prefill again elsewhere.

But locality is not the only objective. Always choosing the warmest worker can create
load imbalance and tail-latency spikes. A router should compare affinity benefit with
queueing cost:
\[
S(w,r)=\alpha\,L_{\rm hit}(w,r)-\beta\,Q(w)-\gamma\,R(w),
\]
where $L_{\rm hit}$ is reusable tokens, $Q$ estimates queue delay, and $R$ represents
resource pressure. This score is illustrative. Real routers need stale-state handling,
bounded coordination overhead, and fairness.

\section{Multi-tenant isolation and timing side channels}

A cache hit changes latency; therefore cache membership can become observable. If
untrusted tenants share a global key space, an attacker can probe candidate prefixes
and statistically infer whether another user recently caused them to be cached.
Published audits have demonstrated this class of leakage in real API systems.

The clean architectural defense is namespace separation. Derive a cache namespace
from the sharing principal---tenant, organization, workspace, or another explicit
trust group---and incorporate it at the root of the identity chain. Requests may
share only when policy deliberately gives them the same namespace.

A raw tenant ID is often not the best external salt. A service can derive
$\kappa=\mathrm{HMAC}_{K_{\rm service}}(\mathrm{principal})$ so internal cache keys
do not expose account identifiers and cannot be chosen adversarially. Key rotation,
model rollout, and cache flush semantics then become explicit operational concerns.

\EngineFigure{ch17-tenant-namespace}{The same token prefix in two security namespaces produces disjoint cache identities; sharing is an authorization decision before it is an optimization.}{fig:prefix-tenants}

\section{Invalidation: model identity is larger than a model name}

A cache created by ``model-x'' yesterday is not automatically compatible with
``model-x'' today. A rollout can change weights, tokenizer files, RoPE scaling,
attention implementation, KV dtype, quantization parameters, adapter weights, or
multimodal preprocessing. Cache identity should use immutable revision/configuration
identifiers, not a mutable marketing alias.

This suggests a practical epoch mechanism. Construct a model-cache epoch from the
artifacts that affect KV semantics. A rollout publishes a new epoch; new requests
cannot hit old entries. Old entries can drain naturally or be reclaimed eagerly.
This turns a subtle invalidation problem into ordinary namespace versioning.

\section{Agents amplify both the benefit and the failure modes}

Agent loops repeatedly send system instructions, tool definitions, conversation
history, and tool results. Without reuse, turn $k$ prefills much of turns
$0\ldots k-1$ again. With exact prefix reuse, the engine computes primarily the newly
appended suffix.

Agents also mutate prompts aggressively. Tool lists may reorder, timestamps may be
injected, retrieval results may be placed before stable instructions, and per-step
IDs may appear near the root. These choices can erase reuse. Prompt architecture and
serving architecture therefore meet at the prefix boundary: stable instructions and
schemas should be canonicalized and placed early when semantics allow, while volatile
data should be appended late.

\section{Observability: prove the cache is helping}

A production dashboard should distinguish lookup activity from useful work avoided.
At minimum record cache-lookup requests, hit tokens, hit-prefix distribution, cold
prefill tokens, recomputed tokens after partial hits, retained bytes, pinned bytes,
eviction count, eviction age, collision-verification failures, stale-handle failures,
and hits partitioned by model epoch and security namespace.

Correlate these with TTFT, queue delay, throughput, and admission failures. A rising
hit rate accompanied by worse p99 latency can indicate that cached blocks are
crowding out active work. A high lookup hit rate with few saved tokens can indicate
pathological tiny prefixes. Metrics need causal interpretation, not celebration.

\section{Failure boundaries}

Four failures deserve explicit tests. First, \emph{false identity}: a collision or
missing namespace field reuses wrong state. Second, \emph{stale identity}: the key is
right but the physical block has been recycled; generation validation must fail.
Third, \emph{unsafe mutability}: two requests share a tail and one appends without
copy-on-write. Fourth, \emph{unauthorized identity}: transformer state would match,
but policy forbids cross-tenant reuse.

The cache should fail closed. A suspicious entry costs recomputation; accepting it can
corrupt model output or leak information.

\EngineFigure{sys-prefix-tree}{An illustrative compressed prefix tree for four requests. Intervals identify logical token spans, not equal token contents across different branches. Two active requests pin their paths; the other leaves have no active readers and are eviction candidates. Actual pin counters and eviction policies are implementation-specific.}{fig:sys-prefix-tree}

\LedgerSection
\begin{ConstraintLedger}
Memory capacity & spends & Retained KV consumes blocks that could admit active requests; only unpinned cache state is reclaimable.\\
Memory bandwidth & trades & Reuse avoids prefill reads/compute but tiered caches may add restore traffic; metadata traversal also touches CPU memory.\\
Compute & buys & A hit avoids forward-pass work for the verified prefix; value scales with reusable tokens and model cost.\\
Latency & trades & Warm prefixes reduce TTFT, while locality-biased scheduling can increase queueing or tail latency.\\
Correctness & risks & Missing dependencies, hash collisions, illegal checkpoints, stale generations, and mutable-tail sharing can all produce wrong state.\\
Cost & trades & More cache memory and coordination can reduce accelerator compute; the optimum depends on workload reuse and admission pressure.\\
\end{ConstraintLedger}

\LabSection{build a prefix cache that can reject its own hits}

The lab in \texttt{code/labs/ch17-prefix-caching/} implements a flat resident-entry
prefix cache with explicit namespaces and a separate prefix-enumeration oracle.
It is not a radix-tree implementation. The tree figures explain the production
representation; this smaller lab isolates equality, isolation and eviction.
It is intentionally
dependency-free. Every lookup verifies exact token ancestry even though the teaching
hash is deliberately weak, making collision behavior testable rather than
astronomically unlikely.

The CLI trace performs insertion, same-tenant reuse and cross-tenant denial.
The tests additionally cover a divergent
suffix, cross-tenant denial, eviction under a block budget, and a manufactured hash
collision, including an exact match hidden behind a colliding candidate.
Tests compare the cache's claimed reusable length against prefix enumeration
over authorized resident entries.

The exercise's implementation contract is concrete: add generation-bearing physical
block handles from Chapter~\ref{ch:paged-kv}; a hit may be committed only if every
handle validates and is pinned. Tests must show that recycled generations are
rejected, pinned blocks cannot be evicted, failed pinning is atomic, and releasing the
last pin makes the block eligible but does not itself require immediate eviction.

\HappenedSection{two industrial shapes, one invariant}

Two influential implementation shapes illustrate the design space. vLLM's current
automatic prefix caching uses hash-derived block identities over ancestry, block
tokens, and extra dependencies. Its documentation describes cache salt for
multi-tenant isolation and cryptographic hash options; current behavior should be
cited to an exact source commit before this chapter is promoted.

SGLang represents prefix topology directly through radix-cache implementations. The
original RadixAttention work couples reuse with the structure of language-model
programs and scheduling opportunities. These systems differ in metadata shape, but
both exploit the same causal fact: only a verified common prefix can reuse exact KV
state.

The historical Hermon claims in the ZERO draft are deliberately not promoted here.
They should be restored only after the exact Hermon source commit and measurement
record are independently rechecked. A textbook should prefer an explicit evidence gap
over a polished but unverified implementation story.

\FrontierSection{2026-10-05}

The original SGLang paper remains the foundational RadixAttention source and reports
up to 6.4$\times$ throughput on its evaluated structured workloads. Current vLLM
documentation shows automatic prefix caching has evolved beyond the early
``one hash per physical block'' simplification: cache identity includes ancestry and
extra dependencies, cryptographic hashing is available/default in current docs, and
matching granularity can be distinct from physical block size in supported
configurations.

Security has also moved from caveat to design constraint. The 2025 paper
\emph{Auditing Prompt Caching in Language Model APIs} demonstrates that cache-dependent
timing can reveal cross-user sharing. The engineering implication is durable even as
individual providers change policy: shared cache namespaces require an explicit trust
decision and threat model.

\ExercisesSection
\begin{enumerate}
\item \emph{Derivation.} A 2,048-token prefix has reuse probability 0.3. Derive
expected FLOPs saved and retained bytes in terms of $P$ and $m_{kv}$. What additional
quantity is needed to compare keeping it with admitting another request?
\item \emph{Trace.} With block size four, trace requests
$[1,2,3,4,5,6]$ and $[1,2,3,4,9]$. Mark reusable full blocks, the divergence, and
which tail must be recomputed.
\item \emph{Correctness.} Give a counterexample showing why
$(\text{token IDs},\text{model name})$ is an insufficient cache key.
\item \emph{Implementation.} Add Chapter~\ref{ch:paged-kv} generation handles to the
lab. Acceptance criteria: stale generations never hit; pin acquisition is atomic;
pinned entries survive eviction pressure; after the final unpin they become eligible.
\item \emph{Security.} Modify the lab so namespaces are derived from a keyed digest
rather than supplied directly. Demonstrate that equal prompts in different principals
cannot share.
\item \emph{Falsification.} Construct a workload where enabling a large prefix cache
reduces goodput because retained blocks prevent admission. State the service objective
under which disabling the cache wins.
\end{enumerate}

\Needspace{18\baselineskip}
\section{Worked problems}

\Needspace{12\baselineskip}
\subsection{Expected value}
\textbf{Problem.} \emph{Derivation.} A 2,048-token prefix has reuse probability 0.3. Derive
expected FLOPs saved and retained bytes in terms of $P$ and $m_{kv}$. What additional
quantity is needed to compare keeping it with admitting another request?

\textbf{Solution.} For length $L=2048$ and reuse probability $h=0.3$,
$E[F_{\rm saved}]=0.3(2P)(2048)=1228.8P$ FLOPs and
$M_{\rm kept}=2048m_{kv}$ bytes. Admission comparison additionally needs the
opportunity cost of those bytes: at minimum the active request's required KV capacity
and service value/queueing consequence.

\Needspace{12\baselineskip}
\subsection{Block trace}
\textbf{Problem.} \emph{Trace.} With block size four, trace requests
$[1,2,3,4,5,6]$ and $[1,2,3,4,9]$. Mark reusable full blocks, the divergence, and
which tail must be recomputed.

\textbf{Solution.} At block size four, both requests have the immutable full block $[1,2,3,4]$.
The first request's tail $[5,6]$ and the second request's tail $[9]$ diverge and are
computed separately. A design publishing only full blocks reuses exactly four tokens.

\Needspace{12\baselineskip}
\subsection{Insufficient key}
\textbf{Problem.} \emph{Correctness.} Give a counterexample showing why
$(\text{token IDs},\text{model name})$ is an insufficient cache key.

\textbf{Solution.} Two deployments may expose the same model name while one has a different LoRA adapter
or weight revision. Equal token IDs then produce different KV. Model aliases are not
immutable computation identities.

\Needspace{12\baselineskip}
\subsection{Falsification}
\textbf{Problem.} \emph{Falsification.} Construct a workload where enabling a large prefix cache
reduces goodput because retained blocks prevent admission. State the service objective
under which disabling the cache wins.

\textbf{Solution.} If retained cold prefixes consume the blocks required to admit high-value arrivals,
queue delay and goodput can worsen despite a high historical hit rate. Under an SLO
that rewards completed requests before deadlines, aggressive eviction or disabling
the cache can win.

